\setglossarystyle{tree}
\makeglossaries

\renewcommand{\glsentryfmt}{\glsentryshort{\glslabel}}

\renewcommand*{\glsgroupskip}{}
\newacronym{acc}{ACC}{uzkrāšana}
\newacronym{ad}{AD}{automatizētā braukšana}
\newacronym{aer}{AER}{adrešu un notikumu reprezentācija}
\newacronym{ai}{AI}{mākslīgais intelekts}
\newacronym{ann}{ANN}{mākslīgais neironu tīkls}
\newacronym{ap}{AP}{vidējā precizitāte}
\newacronym{asic}{ASIC}{lietojumspecifiskā integrālā shēma}
\newacronym{bev}{BEV}{skats no augšas}
\newacronym{bn}{BN}{pakešu normalizācija}
\newacronym{bptt}{BPTT}{kļūdas atpakaļizplatīšana laikā}
\newacronym{bntt}{BNTT}{pakešu normalizācija laikā}
\newacronym{bram}{BRAM}{bloku brīvpiekļuves atmiņa}
\newacronym{cmos}{CMOS}{komplementārā metāla–oksīda–pusvadītāja tehnoloģija}
\newacronym{cnn}{CNN}{konvolucionālais neironu tīkls}
\newacronym{cordic}{CORDIC}{koordinātu rotācijas digitālais skaitļotājs}
\newacronym{clb}{CLB}{konfigurējams loģikas bloks}
\newacronym{dgr}{DGR}{dziļā globālā reģistrācija}
\newacronym{dma}{DMA}{tiešā piekļuve atmiņai}
\newacronym{dnn}{DNN}{dziļais neironu tīkls}
\newacronym{dl}{DL}{dziļā mācīšanās}
\newacronym{dvs}{DVS}{dinamiskās redzes sensors}
\newacronym{dsl}{DSL}{domēnspecifiska valoda}
\newacronym{dsp}{DSP}{digitālā signālu apstrāde}
\newabbreviation{ec}{EC}{evolucionējoša savienojamība}
\newabbreviation{fifo}{FIFO}{pirmais iekšā, pirmais ārā}
\newabbreviation{fpga}{FPGA}{laukā programmējams loģisko vārtu masīvs}
\newabbreviation{ff}{FF}{trigeris}
\newabbreviation{fps}{FPS}{kadri sekundē}
\newabbreviation{fov}{FoV}{redzes lauks}
\newabbreviation{fc}{FC}{pilnībā savienots}
\newabbreviation{gpu}{GPU}{grafiskās apstrādes bloks}
\newabbreviation{gnn}{GNN}{grafu neironu tīkli}
\newabbreviation{hbm}{HBM}{liela joslas platuma atmiņa}
\newabbreviation{hh}{HH}{Hodžkina–Hakslija modelis}
\newabbreviation{if}{IF}{integrējošais impulsu neirons}
\newabbreviation{isi}{ISI}{starpimpulsu intervāls}
\newabbreviation{iob}{IOB}{ieejas un izejas bloki}
\newabbreviation{iou}{IoU}{šķēluma un apvienojuma attiecība}
\newabbreviation{lif}{LIF}{integrējošais neirons ar noplūdi}
\newabbreviation{lut}{LUT}{uzmeklēšanas tabula}
\newabbreviation{ltp}{LTP}{ilgtermiņa potenciācija}
\newabbreviation{ltd}{LTD}{ilgtermiņa depresija}
\newabbreviation{map}{mAP}{klašu vidējā precizitāte}
\newabbreviation{mse}{MSE}{vidējā kvadrātiskā kļūda}
\newabbreviation{mmcm}{MMCM}{jaukta režīma takts vadības bloki}
\newabbreviation{MPSoC}{MPSoC}{daudzprocesoru sistēma mikroshēmā}
\newabbreviation{mpu}{MPU}{mikroprocesora bloks}
\newabbreviation{mac}{MAC}{reizināšana un uzkrāšana}
\newabbreviation{mpbn}{MPBN}{membrānas potenciāla pakešu normalizācija}
\newabbreviation{npu}{NPU}{neirālās apstrādes bloki}
\newabbreviation{nas}{NAS}{neironu arhitektūras meklēšana}
\newabbreviation{nhw}{NHW}{neiromorfā aparatūra}
\newabbreviation{ocm}{OCM}{mikroshēmas iekšējā atmiņa}
\newabbreviation{pe}{PE}{apstrādes elementi}
\newabbreviation{pu}{PU}{apstrādes bloki}
\newabbreviation{pta}{PTA}{konveijerizēts kokveida summators}
\newabbreviation{pll}{PLL}{fāzes automātiskās pieskaņošanas cilpa}
\newabbreviation{pwl}{PWL}{pa daļām lineārs}
\newacronym{ptq}{PTQ}{pēcapmācības kvantēšana}
\newacronym{qat}{QAT}{apmācība, ievērojot kvantēšanu}
\newabbreviation{relu}{ReLU}{taisnotā lineārā vienība}
\newabbreviation{resnet}{ResNet}{atlikuma neironu tīkls}
\newabbreviation{rnn}{RNN}{rekurentie neironu tīkli}
\newabbreviation{roc}{ROC}{rangu secības kodēšana}
\newabbreviation{rsnn}{RSNN}{rekurents impulsu neironu tīkls}
\newabbreviation{ssd}{SSD}{viena posma detektors}
\newabbreviation{stts}{STTS}{telpiski temporālās topoloģijas paraugošana}
\newabbreviation{snn}{SNN}{impulsu neironu tīkls}
\newabbreviation{sdram}{SDRAM}{sinhronā dinamiskā brīvpiekļuves atmiņa}
\newabbreviation{spad}{SPAD}{vienfotona lavīnas fotodiode}
\newabbreviation{stdp}{STDP}{no impulsu laika atkarīgā plastiskuma metode}
\newabbreviation{tdbn}{tdBN}{no sliekšņa atkarīga pakešu normalizācija}
\newabbreviation{tdc}{TDC}{laika–ciparu pārveidotājs}
\newabbreviation{tdp}{TDP}{projektētā siltuma jauda}
\newabbreviation{tvc}{TVC}{temporālā vokseļu kodēšana}
\newabbreviation{tof}{ToF}{lidojuma laiks}
\newabbreviation{tebn}{TEBN}{temporāli efektīvā pakešu normalizācija}
\newacronym{ttfs}{TTFS}{laiks līdz pirmajam impulsam}

